\ifdefined\ishandout
\documentclass[11pt,handout]{beamer}
\else
\documentclass[11pt]{beamer}
\fi

% 日本語対応
\usepackage{luatexja}
\usepackage{luatexja-fontspec}
\renewcommand{\kanjifamilydefault}{\gtdefault}

\usepackage{mathptmx}
\renewcommand{\sfdefault}{lmss}
\renewcommand{\familydefault}{\sfdefault}
\usepackage[T1]{fontenc}
\usepackage{amsmath}
\usepackage{amssymb}
\usepackage{graphicx}
\PassOptionsToPackage{normalem}{ulem}
\usepackage{ulem}
\usepackage{caption}
\captionsetup{labelformat=empty}
\usepackage{bbm}
\usepackage{upgreek}
\setbeamertemplate{section in toc}[sections numbered]
\makeatletter
\usepackage{caption} 
\usepackage{bm}
\usepackage{subfig}
\captionsetup[table]{skip=10pt}
\usepackage{pdflscape}
\newcommand\makebeamertitle{\frame{\maketitle}}%
\AtBeginDocument{%
	\let\origtableofcontents=\tableofcontents
	\def\tableofcontents{\@ifnextchar[{\origtableofcontents}{\gobbletableofcontents}}
	\def\gobbletableofcontents#1{\origtableofcontents}
}

\def\Tiny{\fontsize{7pt}{8pt}\selectfont}
\def\Normal{\fontsize{8pt}{10pt}\selectfont}

\usetheme{Madrid}
\usecolortheme{lily}
\useinnertheme{rounded}

\setbeamertemplate{footline}{\hfill\Normal{\insertframenumber/\inserttotalframenumber}}
\setbeamertemplate{navigation symbols}{}

\newenvironment{changemargin}[2]{%
	\begin{list}{}{%
			\setlength{\topsep}{0pt}%
			\setlength{\leftmargin}{#1}%
			\setlength{\rightmargin}{#2}%
			\setlength{\listparindent}{\parindent}%
			\setlength{\itemindent}{\parindent}%
			\setlength{\parsep}{\parskip}% 
		}%
		\item[]}{\end{list}}

\setbeamertemplate{footline}{\hfill\insertframenumber/\inserttotalframenumber}
\setbeamertemplate{navigation symbols}{}

\usepackage{graphicx}
\usepackage{epsfig}
\usepackage{bm}
\usepackage{epsf}
\usepackage{float}
\usepackage[final]{pdfpages}
\usepackage{multirow}
\usepackage{colortbl}
\usepackage{xkeyval}
\usepackage{listings}
\usepackage{ifthen}
\usepackage{tikz}

\usepackage{xcolor,multirow,colortbl}
\usepackage{verbatim}
\usepackage{setspace}
\usepackage{ragged2e}

\setbeamersize{text margin left=1em,text margin right=1em} 

% Table formatting
\usepackage{booktabs}

% Decimal align
\usepackage{dcolumn}
\newcolumntype{d}[0]{D{.}{.}{5}}

\global\long\def\expec#1{\mathbb{E}\left[#1\right]}
\global\long\def\var#1{\mathrm{Var}\left[#1\right]}
\global\long\def\cov#1{\mathrm{Cov}\left[#1\right]}
\global\long\def\prob#1{\mathrm{Prob}\left[#1\right]}
\global\long\def\one{\mathbf{1}}
\global\long\def\diag{\operatorname{diag}}
\global\long\def\expe#1#2{\mathbb{E}_{#1}\left[#2\right]}
\DeclareMathOperator*{\plim}{\text{plim}}

\usepackage{appendixnumberbeamer}
\renewcommand{\thefootnote}{}

\setbeamertemplate{footline}
{
	\leavevmode%
\hfill
	\insertframenumber{}\hspace*{2ex}\vspace{1ex}
}

\definecolor{blue}{RGB}{0, 0, 210}
\definecolor{red}{RGB}{170, 0, 0}

\makeatother

\usepackage[english]{babel}

\usepackage{tikz}
\newcommand*\circled[1]{\tikz[baseline=(char.base)]{             \node[circle,ball color=structure.fg, shade,   color=white,inner sep=1.2pt] (char) {\tiny #1};}} 

\makeatletter
\let\save@measuring@true\measuring@true
\def\measuring@true{%
	\save@measuring@true
	\def\beamer@sortzero##1{\beamer@ifnextcharospec{\beamer@sortzeroread{##1}}{}}%
	\def\beamer@sortzeroread##1<##2>{}%
	\def\beamer@finalnospec{}%
}
\makeatother

\setbeamersize{text margin left= .8em,text margin right=1em} 
\newenvironment{wideitemize}{\itemize\addtolength{\itemsep}{10pt}}{\enditemize}
\newenvironment{wideitemizeshort}{\itemize}{\enditemize}

\newcommand{\indep}{\perp\!\!\!\!\perp} 

\DeclareMathOperator*{\argmax}{arg\,max}
\DeclareMathOperator*{\argmin}{arg\,min}

\begin{document}
	
	%% タイトルスライド
	\begin{frame}[noframenumbering]{}
		\vspace{0.5cm}
		\title[]{復習セッション}
		\author{Jonathan Roth}
		\date{数理計量経済学 I \\ ブラウン大学\\} 
		\titlepage {\small{}\ }\thispagestyle{empty} \vspace{-30pt}
		
	\end{frame}


\begin{frame}{概要}
	\begin{enumerate}
		\item 
		期末試験の構成と運営について
		\vspace{0.5cm}
		\item
		講義の主要な概念
		\vspace{0.5cm}
		\item
		質疑応答
	\end{enumerate}
\end{frame}
	
	
\begin{frame}{期末試験の構成}
	\begin{wideitemize}
		\item
		構成は昨年の期末試験と同様です。
		
		\item
		このコースで学んだ数学的ツールを使用する分析的な問題が1〜2問出題されます。
		
		\item
		新しい実証研究の応用例に関する問題が出題されます。因果的効果を学ぶための仮定を述べ、評価し、実証結果についてコメントすることが求められます。
		
	\end{wideitemize}
\end{frame}	


\begin{frame}{期末試験の運営について}
	\begin{wideitemize}
		
		\item
		試験は対面形式で、時間は2時間です。
			\begin{wideitemize}
				\item
				追加時間の延長などの合理的配慮（SAS accommodation）が必要な場合は、私とTAにメールをしてください。
			\end{wideitemize}
		
		\item
		日時は12月15日午後2時、場所は MacMillan Hall 117 です。
		
		\item
		A4サイズ（8.5x11インチ）2枚のメモ（自筆・印刷不問）を持ち込むことができます。
			\begin{itemize}
				\item 
				例：中間試験用の「チートシート」に加えて、新しいものを1枚。
			\end{itemize}
		
	\end{wideitemize}
\end{frame}

\begin{frame}{コースの概要}
	\begin{wideitemizeshort}
		\item
		このクラスでは、主にデータを用いて経済学的な\underline{因果的}な問いにどのように答えるかに焦点を当てました。
			\begin{itemize}
				\item
				ブラウン大学とURIのどちらに進学するかが、所得に与える影響は？ 
				
				\item
				健康保険への加入が、うつ病に与える影響は？ 
				
				\item
				最低賃金が雇用に与える影響は？
			\end{itemize}
		
		\pause
		\item
		因果制の効果の考え方を、\textit{潜在的結果}（potential outcomes）を用いて定式化しました。
			\begin{itemize}
				\item 
				各ユニットは、処置を受けた場合と受けなかった場合の両方の潜在的結果 $Y_i(1)$ と $Y_i(0)$ を持ちます。
				
				\item
				例： $Y_i(1) =$ ブラウン大に行った場合の所得、 $Y_i(0) =$ URIに行った場合の所得。
				
				\item
				観測できるのは、処置群については $Y_i(1)$、対照群については $Y_i(0)$ のみです。
				
				\item
				関心があるのは、因果的効果 $Y_i(1) - Y_i(0)$ （またはその母集団平均）です。
			\end{itemize}
	\end{wideitemizeshort}
\end{frame}

\begin{frame}{2つの主要な課題}

データで因果的な問いに答えるには、2つの主要な課題があります： \medskip 
	
	\begin{wideitemize}
		\item
		各ユニットについて、反実仮想的な結果（counterfactual outcome）を観察することは決してできません。
			\begin{itemize}
				\item
				例：ブラウン大生については所得 ($Y_i(1)$) を観察できますが、もし彼らがURIに行っていた場合の所得 ($Y_i(0)$) は分かりません。
			\end{itemize} 
		
		\item
		通常、私たちが関心を持つ母集団全体ではなく、そのサンプル（標本）のデータしか観察できません。
		\begin{itemize}
			\item
			例：最近の卒業生の一部の調査データしか持っていません。
		\end{itemize}
	\end{wideitemize}

\end{frame}

\begin{frame}{識別 vs 統計的推論}

通常、これら2つの問題を分けて考えます： \medskip	
	
	\pause 
	
	\begin{wideitemize}
		\item
		\textbf{識別（Identification）：} もし母集団全体の観察可能なデータが手に入るとしたら、因果的効果について何を学ぶことができるか？ 
		\pause
			\begin{wideitemize}
				\item
				通常、処置がどのように割り当てられるかについての仮定から始めます。
				
				\item
				それらの仮定の下で、因果的効果が観察可能な母集団平均の関数であることを示します。
			\end{wideitemize}
		\pause
		\item
		\textbf{統計的推論（Statistical Inference）：} サンプルが与えられたとき、母集団の観察可能な特徴について何を学ぶことができるか？ \pause
			\begin{wideitemize}
				\item
				通常、母集団平均を標本平均で推定（代入）します。
				
				\item
				条件付き平均を推定する必要があるときは、回帰を用いて条件付き期待値関数（CEF）を近似します。
				
				\item
				統計的ツールを用いて、仮説検定を行ったり信頼区間を構築したりします。
			\end{wideitemize}		
	\end{wideitemize}
\end{frame}

\begin{frame}{5種類の識別論拠}
	\begin{wideitemize}
		\item
		\textbf{実験（Experiments）：} 処置がランダム化されていれば、処置群と対照群の結果を直接比較できます。
		
		\pause
		\item
		\textbf{条件付き非交絡性（Conditional unconfoundedness）：} 観察可能な特徴を条件とすれば処置の割り当てが実験のようであると仮定し、同じ共変量を持つ処置群と対照群を比較します。
		
		\pause
		\item
		\textbf{差の差分析（Difference-in-differences）：} 処置への選択を許容しますが、選択バイアスが時間を通じて一定であると仮定し、処置前後の変化の差を比較します。
		
		\pause
		\item
		\textbf{操作変数法（Instrumental variables）：} 操作変数の割り当てがランダムであり、処置を通じてのみ結果に影響を与えると仮定し、操作変数の値が異なる集団間を比較します。
		
		\pause
		\item
		\textbf{回帰不連続デザイン（Regression discontinuity）：} カットオフの周囲で交絡因子が連続的に変化すると仮定し、スコアがカットオフの直下の人々と直上の人々を比較します。
	\end{wideitemize}	
\end{frame}

\begin{frame}{実験}

うつ病に関する標本平均
\begin{tabular}{lll}
	& 対照群 & 処置群 \\
	平均 & 0.329 & 0.306\\
	標準偏差 (SD) & 0.470 &  0.461 \\
	サンプルサイズ (N) & 10426 & 13315  
\end{tabular}


\begin{wideitemize}
	\item
	識別： \pause{}$D_i \indep (Y_i(1),Y_i(0))  \implies ATE = E[Y_i | D_i =1 ] - E[Y_i | D_i = 0]$ \\ \medskip これをどのように導出するか確認しておいてください！
	
	\pause
	\item
	推定：母平均を標本平均に置き換えます！
	
	\pause
	\item
	推定（その2）：OLS でも推定可能です。
	$$Y_i = \alpha + \beta D_i + \epsilon_i$$
\end{wideitemize}
	
\end{frame}


\begin{frame}{条件付き非交絡性}
	\begin{wideitemize}
		\item
		例：どの大学に行くかは、合格した大学のセットを条件とすれば実質的にランダムかもしれません（Dale \& Krueger）。
		
		\pause
		\item
		識別： $D_i \indep (Y_i(1),Y_i(0)) | X_i  \implies CATE(x) = E[Y_i | D_i =1 , X_i = x] - E[Y_i | D_i = 0, X_i = x]$ 
		
		\pause
		\item
		推定：通常、CEF を近似する必要があります $\rightarrow$ OLS を使用します！
		
		\item
		以下の OLS 推定値を用いて ATE を近似するのが一般的です：
		
		$$Y_i = \alpha + \beta D_i + \bm{\gamma}' \bm{X_i} + \epsilon_i$$
		
		\pause
		\item
		これは、(i) 条件付き非交絡性が成り立ち、かつ (ii) CEF が近似的に線形、すなわち $E[Y_i | D_i, \bm{X}_i] \approx \alpha + \beta D_i + \bm{\gamma}' \bm{X_i} $ である場合にうまく機能します。
	\end{wideitemize}	
\end{frame}

\begin{frame}
	\centering
	\includegraphics[width = 0.5\linewidth]{dk-results-table-reg2}
\end{frame}

\begin{frame}{差の差分析 I}
	\centering
	\includegraphics[width = 0.65\linewidth]{../Chapter6/hastings-event-study-w-labels}
	
	\begin{wideitemize}
		\item
		主要な識別仮定：並行トレンド --- 選択バイアスが時間を通じて一定であること（形式的な定義を確認しておいてください）。
		
		\pause
		\item
		識別：並行トレンド（および事前反応なし）の下で、 $\tau_{ATT} = \underbrace{(\mu_{12} - \mu_{11})}_{\text{処置群の変化}} - \underbrace{ (\mu_{02} - \mu_{01}) }_{\text{対照群の変化}}$
	
		\pause	
		\item
		推定：母平均の代わりに標本平均を代入します！
		
		\pause
		\item
		推定（その2）：OLS でも推定可能です（方法を確認しておいてください！）。
	\end{wideitemize}
\end{frame}

\begin{frame}{差の差分析 II}
	\begin{wideitemize}
		\item
		DID の仮定の妥当性を評価するために、「イベントスタディ」を用いて処置前のトレンドを確認するのが一般的です。
		
		\centering
		\includegraphics[width = 0.5\linewidth]{../Chapter6/medicaid-eligibility}
		
		\item
		(i) プレトレンドが 0 に近く、かつ (ii) すべての信頼区間（CI）を通る直線が引けない（トレンドに変化がある）場合に、リサーチデザインに対する信頼が高まります。
		
		\pause
		\item
		イベントスタディの推定も OLS で行えます。
	\end{wideitemize}
\end{frame}

\begin{frame}{操作変数法 (IV)}
	\centering
	\includegraphics[width = 0.6\linewidth]{../Chapter7/ak-first-stage-basic}
	
	\begin{wideitemize}
		\item
		IV では、実質的にランダムに割り当てられ、処置への影響を通じてのみ結果に影響を与える「操作変数」を用います。
		
		\item
		4つの主要な識別仮定（これらを理解しておいてください！）：
		\begin{wideitemizeshort}
			\item
			独立性：操作変数は実質的にランダムに割り当てられている。
			
			\item
			除外制約：操作変数は処置を通じてのみ結果に影響を与える。
			
			\item
			単調性：ディファイアーが存在しない。
			
			\item
			関連性：操作変数が処置ステータスに影響を与える。
		\end{wideitemizeshort}	
		 
		
		\pause
		\item
		これら4つの仮定の下で、 $\beta_{IV} = \frac{ E[Y_i | Z_i = 1] - E[Y_i | Z_i = 0] }{E[D_i | Z_i = 1] - E[D_i | Z_i = 0]}$ は LATE（コンプライヤーにおける平均処置効果）を識別します。
	\end{wideitemize}
\end{frame}

\begin{frame}{操作変数法 --- 推定}
	\centering
	\includegraphics[width = 0.6\linewidth]{../Chapter7/ak-first-stage-basic}

	\begin{wideitemize}
		\item
		これら4つの仮定の下で、 $\beta_{IV} = \frac{ E[Y_i | Z_i = 1] - E[Y_i | Z_i = 0] }{E[D_i | Z_i = 1] - E[D_i | Z_i = 0]}$ は LATE を識別します。
		
		\item
		推定は、誘導型と第一段階に標本対応物を代入し、その比をとることで行われます。
		
		\pause
		\item
		推定（その2）：「二段階最小二乗法（2SLS）」でも実行可能です。これにより複数の操作変数を取り込むことができます（方法を確認しておいてください）。
			\begin{wideitemize}
				\item
				$D_i$ を操作変数 $Z_i$ （およびコントロール変数）に回帰する。
				
				\item
				$Y_i$ をその予測値 $\hat{D}_i$ （およびコントロール変数）に回帰する。
			\end{wideitemize}
	\end{wideitemize}
	
\end{frame}

\begin{frame}{回帰不連続デザイン (RDD)}
	\centering
	\includegraphics[width = 0.6 \linewidth]{../Chapter8/bleemer-rf}
	
	\begin{wideitemize}
		\item
		RDD では、処置ステータスを（部分的に）決定するしきい値の直上と直下の人々を比較します。
		
		\pause
		\item
		識別：潜在的結果がカットオフにおいて連続であること。
			\begin{itemize}
				\item 
				これがどのような場合に満たされないか（操作など）を理解しておいてください！
			\end{itemize}
		
		\pause
		\item
		推定：OLS または局所線形回帰を用いて、カットオフにおける CEF を推定します。
			\begin{wideitemize}
				\item
				注意：局所線形回帰の詳細については試験には出しません。
			\end{wideitemize}
	\end{wideitemize}
\end{frame}


\begin{frame}{おわりに I}
この1学期で多くのことを学びました！ \medskip

\begin{wideitemize}
	\item
	因果的効果を推定することの難しさを学びました。
	
	\item
	いつ、どのように因果的効果について学べるかを考えるための統計的な言語を習得しました。
	
	\item
	因果的効果を学ぶための、いくつかの実践的な「識別戦略」を学びました。
	
	\item
	有限のサンプルにおいて（多くの場合、回帰を用いて）因果的効果を推定し、仮説検定を行うためのツールを習得しました。
\end{wideitemize}

\pause
\medskip
さらに深く学びたい方は、他にも多くのクラスがあります！
	\begin{itemize}
		\item 
		ノンパラメトリック手法、機械学習アプローチ、ベイズ計量経済学、時系列分析と予測、など。
	\end{itemize}

\end{frame}

\begin{frame}{おわりに II}
これらのツールを、今後さまざまな形で応用してほしいと願っています： \medskip

\begin{wideitemize}
	\item
	政策や企業の意思決定を改善し、社会福祉（あるいは企業の利益 --- お好みで！）を向上させるための研究を行う。
	
		\begin{itemize}
			\item 
			学術研究に興味があるなら、ブラウン大の教授のリサーチ・アシスタント（RA）をしたり、卒業論文を書いたりすることをお勧めします。
		\end{itemize}
	
	\item
	新聞やオンラインの記事などを読む際に、実証的な証拠をより深く理解する。
	
	\item
	計量経済学者になって、将来のより良い政策分析のためのツールを開発する :-) 
	
	
\end{wideitemize}

	
\end{frame}


\begin{frame}{おわりに III}
	\begin{wideitemize}
		\item
		コース評価（フィードバック・フォーム）への記入をお願いします： \href{https://brown.evaluationkit.com  }{\underline{https://brown.evaluationkit.com}}
		
		\item
		皆さんからのフィードバックは、評価のためだけでなく、今後のコース改善のために活用されます！
		
		\item
		私がこのコースを教えるのは3回目ですが、何がうまくいき、何が改善できるかについてのコメントをいただければ幸いです。ありがとうございました！
	\end{wideitemize}
\end{frame}

\begin{frame}
	\centering
	\includegraphics[width = 0.9 \linewidth]{poppy-good-luck}
\end{frame}

\end{document}
